\section{Exact localization and removal of amplification}
\label{sec:localization}

We now use Proposition~\ref{prop:perturbation} to prove the fixed-amplification
part of Theorem~\ref{thm:exact-infiniteness}.  Throughout this section,
let $C$ be an exact $C^*$-algebra satisfying \propF{} for one fixed
integer $n\geq1$, and retain the constant $0<\sigma<1$ from
\eqref{eq:perturbation-constant}.  For $P\in\Prim(C)$, write
$a_P=a+P$, with its norm computed in any irreducible representation
having kernel $P$.  A hereditary subalgebra $\mathcal D\subset C$
is nonzero at $P$ when $\mathcal D\not\subset P$.

Our immediate objective is to construct a square-zero element in a
hereditary subalgebra that is nonzero on a prescribed quasi-compact
subset of $\Prim(C)$.  Its initial and range squares are orthogonal
and Cuntz equivalent.  Repeating the construction on one fixed norm
slice will remove the amplification integer $n$ and give \propP{}.
Theorem~\ref{thm:individual-to-strong} will then supply the strong
pure infiniteness conclusion.  All finite covers in the construction
concern the primitive ideal space; the partition of unity used below
is on a real interval.

\subsection{Local elements and finite covers}

The next lemma uses only \propF{}: finite-dimensional hereditary
images are excluded, and irreducibility then supplies common positive
elements locally.  Exactness enters later through
Proposition~\ref{prop:perturbation}, when these local choices must
work simultaneously on a whole quasi-compact set.

\begin{lemma}[Local common elements]\label{lem:local-common-element}
Let $C$ be an exact $C^*$-algebra satisfying \propF{} for one fixed
integer $n\geq1$, and let $P\in\Prim(C)$.
\begin{enumerate}[label=\textup{(\roman*)}]
\item If $h\in C_+$ and $h_P\ne0$, then for every integer $m\ge1$
there are $m$ pairwise orthogonal elements of $\Her(h)_+$, each
nonzero at $P$.
\item If $h,k\in C_+$ are nonzero at $P$, there is $c\in C_+$
such that $c_P\ne0$ and
\[
 c^{\oplus n}\precsim h,\qquad c^{\oplus n}\precsim k.
\]
If $h,k$ belong to a hereditary subalgebra $\mathcal D$, $c$ can also
be chosen in $\mathcal D$.
\end{enumerate}
\end{lemma}

\begin{proof}
Fix an irreducible representation $\pi:C\to B(H_\pi)$ with kernel
$P$.  First, no nonzero $\pi(a)$, $a\in C_+$, has finite rank.
Indeed, if its rank were $d\in\N$, applying the amplified
representation to
$a^{\oplus 2n}\precsim a^{\oplus n}$ would express an operator of
rank $2nd$ as a norm limit of operators of rank at most $nd$.
This is impossible, since the set of operators of rank at most a
fixed finite integer is norm closed.

Put $H_h=\overline{\pi(h)H_\pi}$.  We claim that the restriction of
$\pi(\Her(h))$ to $H_h$ is strongly dense in $B(H_h)$.
Irreducibility and the bicommutant theorem
\cite[I.9.1.1]{Blackadar} give strong density of $\pi(C)$ in $B(H_\pi)$.  Consequently the strong closure of
$\pi(\Her(h))$ contains $\pi(h)B(H_\pi)\pi(h)$.
For $t>0$, let $E_t=1_{[t,\infty)}(\pi(h))$.  Every operator
$E_tTE_t$, with $T\in B(H_h)$ extended by zero on $H_h^\perp$, has
the form $\pi(h)S\pi(h)$: the inverse of $\pi(h)$ on $E_tH_\pi$
is bounded.  As $t\downarrow0$, $E_tTE_t\to T$ strongly on $H_h$.
This proves the claim.

The space $H_h$ is infinite dimensional by the finite-rank
obstruction.  Choose orthonormal vectors $\xi_1,\ldots,\xi_m$ in
$H_h$ and distinct numbers $\lambda_1,\ldots,\lambda_m\in(0,1)$.
Let $T\in B(H_h)$ be the positive contraction with
$T\xi_j=\lambda_j\xi_j$ and zero on their orthogonal complement.
Choose nonnegative continuous functions $f_j$ on $\R$, with disjoint
compact supports in $(0,1)$ and $f_j(\lambda_j)=1$.
Kaplansky density \cite[I.9.1.3]{Blackadar} supplies a bounded net of selfadjoint elements
$b_\lambda$ of the represented hereditary algebra converging
strongly to $T$.  Continuous functional calculus on bounded
selfadjoint sets is strongly continuous; this follows by uniform
polynomial approximation on a common bounded interval.  Hence
$f_j(b_\lambda)\xi_j\to\xi_j$.  One $b_\lambda$ makes all these
$m$ vectors nonzero.  Lift it to a selfadjoint
$b\in\Her(h)$.  The elements $f_j(b)$ lie in $\Her(h)$, are
pairwise orthogonal, and are nonzero at $P$.  The functions vanish
at zero, so this argument also applies to a nonunital hereditary
algebra.  This proves \textup{(i)}.

For two positives $u,v$ nonzero at $P$, strong density of $\pi(C)$
gives $x\in C$ with $\pi(uxv)\ne0$.  For example, choose a vector
on which $\pi(v)$ is nonzero and a rank-one operator carrying its
image to a vector not annihilated by $\pi(u)$; approximate that
operator strongly on the image under $\pi(v)$.  With $z=uxv$, the positive
$d=zz^*$ satisfies
\[
 d_P\ne0,\qquad d\in\Her(u),\qquad
 d\CuEq z^*z\in\Her(v).
\]
Here and below $\CuEq$ denotes Cuntz equivalence, and the comparison
rules are those of Lemma~\ref{lem:comparison-rules}.
For a finite list of positives nonzero at $P$, repeating this
operation produces one positive, nonzero at $P$, subequivalent to
every element of the list.  At each step the new element lies in
the hereditary algebra of the preceding one, so all previous
comparisons and any prescribed hereditary membership persist.

Apply \textup{(i)} with $m=n$ separately to $h$ and $k$, obtaining
orthogonal lists $h_1,\ldots,h_n$ and $k_1,\ldots,k_n$.
Take a common element $c$ for the resulting list of $2n$ positives.
Then
\[
 c^{\oplus n}\precsim\bigoplus_{j=1}^n h_j
 \CuEq\sum_{j=1}^n h_j\precsim h,
\]
and the same argument gives $c^{\oplus n}\precsim k$.
The asserted choice in $\mathcal D$ follows from the hereditary
membership just observed.
\end{proof}

We record the norm-attainment fact that ensures a strict inequality
in the finite-cover argument.

\begin{lemma}[Pure-state norm attainment]\label{lem:purestateattainment}
If $B$ is a $C^*$-algebra and $0\ne b\in B_+$, there is an
irreducible representation $\pi$ of $B$ with
$\|\pi(b)\|=\|b\|$.
\end{lemma}

\begin{proof}
In the state space of the unitization $\widetilde B$, the set
\[
 F_b=\{\varphi:\varphi(b)=\|b\|\}
\]
is nonempty: evaluation at the largest spectral value of $b$ on
$C^*(1,b)$ extends to a state of $\widetilde B$ by the positive
Hahn--Banach extension theorem \cite[II.6.3.1]{Blackadar}.  It is a weak-star compact face,
since every state has value at most $\|b\|$ on $b$.  By the
Krein--Milman theorem it has an extreme point, which is therefore
a pure state of $\widetilde B$; see \cite[II.6.3.2--3]{Blackadar}.
Let $\widetilde\pi$ be its irreducible GNS representation
\cite[II.6.4.8]{Blackadar}.  It does not annihilate $B$, because
its value on $b$ is positive.  The closure of
$\widetilde\pi(B)H_{\widetilde\pi}$ is a nonzero reducing subspace
for $\widetilde\pi(\widetilde B)$, hence the whole space.
If an operator commutes with $\widetilde\pi(B)$, then its commutator
with $\widetilde\pi(x)$ annihilates $\widetilde\pi(B)H_{\widetilde\pi}$
for every $x\in\widetilde B$, since $xB\subset B$.
It therefore commutes with $\widetilde\pi(\widetilde B)$.
Thus the restriction $\pi=\widetilde\pi|_B$ is nondegenerate and
irreducible.  The vector state gives
$\|\pi(b)\|\ge\|b\|$, and contractivity gives equality.
\end{proof}

\begin{lemma}[Primitive norm slices]\label{lem:primitive-compactness}
Let $C$ be a $C^*$-algebra.  For $c\in C_+$ and $\alpha>0$, the set
\[
 \Omega(c,\alpha)=\{P\in\Prim(C):\|c_P\|>\alpha\}
\]
is open, and
\[
 K(c,\alpha)=\{P\in\Prim(C):\|c_P\|\ge\alpha\}
\]
is quasi-compact.  No Hausdorff hypothesis is required.
\end{lemma}

\begin{proof}
The first set is $\{P:(c-\alpha)_+\notin P\}$, which is open in
the hull-kernel topology.  For the second assertion, extend any
relative open cover to ambient open sets.  Write those sets as
\[
 \mathcal O_\lambda=\{P:J_\lambda\not\subset P\},
 \qquad I=\overline{\sum_\lambda J_\lambda},
\]
where each $J_\lambda$ is a closed ideal of $C$.
If $\|c+I\|\ge\alpha$, Lemma~\ref{lem:purestateattainment}
applied in $C/I$ gives an irreducible representation attaining that
norm.  Its pullback has a primitive kernel containing $I$, belongs
to $K(c,\alpha)$, and lies outside every $\mathcal O_\lambda$.
This contradicts the covering property.  Thus
$\|c+I\|<\alpha$; if $C/I=0$, the same inequality is immediate.
Choose $d\in I$ with $\|c-d\|<\alpha$, and then approximate $d$
by an element
$d_F\in\sum_{\lambda\in F}J_\lambda$ for a finite index set
$F$, closely enough that $\|c-d_F\|<\alpha$ still holds.
Any primitive ideal containing every $J_\lambda$, $\lambda\in F$,
has $\|c_P\|<\alpha$.  Hence the finitely many
$\mathcal O_\lambda$, $\lambda\in F$, cover $K(c,\alpha)$.
The proof uses quasi-compactness only and does not assert that
$K(c,\alpha)$ is closed.
\end{proof}

\begin{lemma}[Finite coupling]\label{lem:finite-coupling}
Let $C$ be an exact $C^*$-algebra satisfying \propF{} for one fixed
integer $n\geq1$.  Let $K_0\subset\Prim(C)$ be quasi-compact and let
$\mathcal D\subset C$ be hereditary.  Suppose that at each
$P\in K_0$ one can choose positive contractions $p,q\in\mathcal D$
and $h,k\in\mathcal D_+$ satisfying
\[
 pq=0,\qquad ph=h,\qquad qk=k,\qquad h_P\ne0,\quad k_P\ne0.
\]
There are finitely many such choices $(p_i,q_i,h_i,k_i)$ and one
selfadjoint $U\in\mathcal D$ such that, at every $P\in K_0$,
\[
 (p_iUq_i)_P\ne0
\]
for at least one selected index $i$.  Any support conditions in
the original choices are retained in the selected ones.
\end{lemma}

\begin{proof}
If $K_0$ is empty, choose $U=0$ and the empty family.  Otherwise,
Lemma~\ref{lem:local-common-element} associates to every original
choice a positive $c$ nonzero at its specified point, with
$c^{\oplus n}\precsim h,k$.  Choose
$0<\alpha<\|c_P\|$.  The resulting open sets
$\Omega(c,\alpha)$ cover $K_0$, so choose a finite subcover
$\Omega_i=\Omega(c_i,\alpha_i)$, $1\le i\le N$.

Set $\theta=\sigma/4$, $\eta_i=\theta^{i-1}$, and $U_0=0$.
At stage $i$, apply Proposition~\ref{prop:perturbation} with
$T=p_iU_{i-1}q_i$, scale $\eta_i$, and the selected data.
The identities $p_ih_i=h_i$, $q_ik_i=k_i$ imply
$p_i^2h_i=h_i$, $q_i^2k_i=k_i$.  The simultaneous hereditary and
selfadjoint options give a selfadjoint contraction
$y_i\in\mathcal D$ such that $U_i=U_{i-1}+\eta_i y_i$ satisfies
\[
 \|(p_iU_iq_i)_P\|\ge\sigma\eta_i
 \qquad(P\in K(c_i,\alpha_i)).
\]
Let $U=U_N$.  Later perturbations change the $i$th corner by at
most
\[
 \sum_{j>i}\eta_j
 \le\eta_i\frac{\theta}{1-\theta}
 =\eta_i\frac{\sigma}{4-\sigma}
 <\frac{\sigma\eta_i}{2}.
\]
Consequently
\begin{equation}\label{eq:finite-coupling-bound}
 \|(p_iUq_i)_P\|>\frac{\sigma\eta_i}{2}
 \qquad(P\in K(c_i,\alpha_i)).
\end{equation}
In particular each selected corner is nonzero throughout its
entire open set $\Omega_i$, proving the assertion.
\end{proof}

When applying Lemma~\ref{lem:finite-coupling}, it suffices to
produce two orthogonal positives $u,v$ nonzero at the specified
point.  Indeed, choose $0<\beta<\|u_P\|$ and a continuous
$f:[0,\infty)\to[0,1]$ that is zero near zero and one on
$[\beta,\infty)$.  Then $p=f(u)$ and $h=(u-\beta)_+$ satisfy
$ph=h$ and $h_P\ne0$.  The analogous construction from $v$ gives
$q,k$.  All four elements retain the hereditary and spectral
support restrictions of $u,v$.

\subsection{A square-zero element on a quasi-compact set}

\begin{proposition}\label{prop:square-zero}
Let $C$ be an exact $C^*$-algebra satisfying \propF{} for one fixed
integer $n\geq1$.  Let $K_0\subset\Prim(C)$ be quasi-compact, and let
$\mathcal D\subset C$ be a hereditary subalgebra nonzero at every
point of $K_0$.  There exists $x\in\mathcal D$ such that
\[
 x^2=0,\qquad x_P\ne0\quad(P\in K_0).
\]
\end{proposition}

\begin{proof}
If $K_0=\varnothing$, take $x=0$.  Assume henceforth that $K_0$
is nonempty.  We first construct a selfadjoint element that is
nonscalar on every relevant hereditary support space.  Coupling its
separated spectral windows and deleting the nearby spectral blocks
will produce a selfadjoint element whose positive and negative parts
are both nonzero at every point of $K_0$.  A final coupling of those
two orthogonal parts gives the required square-zero element.

\paragraph{First coupling: a nonscalar selfadjoint element.}
At each $P\in K_0$, choose $h\in\mathcal D_+$ with $h_P\ne0$.
Lemma~\ref{lem:local-common-element}\textup{(i)}, with $m=2$,
provides two orthogonal positives in $\Her(h)\subset\mathcal D$
nonzero at $P$.  Apply the preceding cutoff construction to these
two positives and use the resulting four inputs in
Lemma~\ref{lem:finite-coupling} to obtain $U=U^*\in\mathcal D$.
For an irreducible representation $\pi$ with kernel $P$, put
$H_{\mathcal D,\pi}=\overline{\pi(\mathcal D)H_\pi}$.
This space reduces $\pi(\mathcal D)$, which acts as zero on its
orthogonal complement.  At $P$ one selected $pUq$ is nonzero.
Since $p,q\in\mathcal D$ and $pq=0$, this implies
\begin{equation}\label{eq:nonscalarU}
 \pi(U)|_{H_{\mathcal D,\pi}}\text{ is not a scalar operator}
 \qquad(P\in K_0).
\end{equation}
In particular, nonscalarity is obtained on the support space of
$\mathcal D$, without using its zero action on the complement.

\paragraph{Second coupling: separated spectral windows.}
At a specified $P\in K_0$, the selfadjoint operator in
\eqref{eq:nonscalarU} has two distinct spectral values.  Choose
disjoint closed intervals $I,J\subset\R$ at positive distance,
and continuous nonnegative functions $F,G$ supported in $I,J$,
respectively, such that both
$F(\pi(U))|_{H_{\mathcal D,\pi}}$ and
$G(\pi(U))|_{H_{\mathcal D,\pi}}$ are nonzero.
The functions may be chosen on all of $\R$ with compact support.
If a window contains zero, its function need not vanish at zero.
We evaluate scalar functions in $\mathcal D+\C1$ inside a fixed
unitization $\widetilde C$ of $C$, and take spectral projections in
$\widetilde C^{**}$.
Let $(d_\lambda)$ be a positive contractive approximate identity of
$\mathcal D$.
The positive elements
\[
 F(U)d_\lambda F(U),\qquad G(U)d_\lambda G(U)
\]
belong to $\mathcal D$.  Their images on
$H_{\mathcal D,\pi}$ converge strongly to $F(\pi(U))^2$ and
$G(\pi(U))^2$, respectively, so a choice of middle elements makes
both images nonzero.  These positives are supported, globally in
the bidual, on $1_I(U)$ and $1_J(U)$.

The same cutoff construction gives the corresponding $p,q,h,k$,
retaining these spectral supports.  Apply
Lemma~\ref{lem:finite-coupling} again to obtain
$V=V^*\in\mathcal D$ and finitely many selected pairs
$(p_i,q_i)$, supported in windows $(I_i,J_i)$, such that
\begin{equation}\label{eq:spectralwindowV}
 \begin{gathered}
 p_i=1_{I_i}(U)p_i1_{I_i}(U),\qquad
 q_i=1_{J_i}(U)q_i1_{J_i}(U),\\
 \text{for every }P\in K_0\text{ there is }i
 \text{ with }(p_iVq_i)_P\ne0.
 \end{gathered}
\end{equation}
Put $\rho=\min_i\dist(I_i,J_i)>0$.
The positive elements $p_i,q_i$ need not commute with $U$;
only the displayed support identities will be used.

\paragraph{Spectral deletion produces both signs.}
We retain the corners in \eqref{eq:spectralwindowV} while removing
all sufficiently narrow diagonal spectral blocks.  A positive operator
with all those diagonal blocks zero must vanish; the same holds for
a negative operator.  Since the retained corners prevent vanishing,
this will force both signs in every quotient over $K_0$.

Fix a compact interval $[-M,M]$ containing the spectrum of $U$
in the unitization.  Take a finite continuous nonnegative
partition of unity $f_1,\ldots,f_s$ on this interval, whose
supports $S_j$ have diameter less than $\rho/8$.
Define
\begin{equation}\label{eq:deletionH}
 H=\sum_{\dist(S_j,S_k)>\rho/2}
       f_j(U)Vf_k(U).
\end{equation}
The summation uses ordered pairs, and its condition is symmetric
in $j,k$, so $H=H^*$.  Although a factor $f_j(U)$ may have a
scalar part, it lies in $\mathcal D+\C1$; the middle
factor $V\in\mathcal D$ implies $H\in\mathcal D$.

For each selected pair $(p_i,q_i)$ we have
\begin{equation}\label{eq:window-preservation}
 p_iHq_i=p_iVq_i.
\end{equation}
Indeed, expand $p_iVq_i$ using $\sum_jf_j(U)=1$ on both sides.
For a summand $p_if_j(U)Vf_k(U)q_i$ to be nonzero, the supports
$S_j,S_k$ must meet $I_i,J_i$, respectively, by
\eqref{eq:spectralwindowV}.  Their distances then satisfy
\[
 \dist(S_j,S_k)
 \ge\dist(I_i,J_i)
       -\operatorname{diam}(S_j)-\operatorname{diam}(S_k)
 >\frac{3\rho}{4}>\frac\rho2.
\]
Every potentially nonzero summand is therefore retained in
\eqref{eq:deletionH}.  This argument requires no commutation of
$p_i,q_i$ with $U$.  Equations~\eqref{eq:spectralwindowV} and
\eqref{eq:window-preservation} show that $H_P\ne0$ for every
$P\in K_0$.

Let $B$ be any Borel interval of length less than $\rho/4$.
The two supports in a retained pair cannot both meet $B$, since
their distance exceeds $\rho/2$.  Hence
\[
 1_B(U)H1_B(U)=0.
\]
Partition $[-M,M]$ into finitely many disjoint Borel intervals
$B_j$ of length less than $\rho/4$, assigning endpoints so that
their spectral projections sum to the identity.  In any unital
extension of a representation, the resulting projections $E_j$
thus satisfy $E_j\pi(H)E_j=0$ and $\sum_jE_j=1$.
If $\pi(H)\ge0$, then
\[
 \|\pi(H)^{1/2}E_j\xi\|^2
 =\langle E_j\pi(H)E_j\xi,\xi\rangle=0
\]
for every $j,\xi$, whence $\pi(H)=0$.  Applying the same argument
to $-\pi(H)$ excludes a nonzero negative semidefinite image.
The spectral projection calculations can equivalently be made
in the bidual and passed through the normal extensions of the
representations.  We conclude that
\begin{equation}\label{eq:bothsigns}
 (H_+)_P\ne0,\qquad (H_-)_P\ne0
 \qquad(P\in K_0).
\end{equation}
Both positive parts belong to $\mathcal D$ and are orthogonal.

\paragraph{Third coupling: the square-zero product.}
At every $P\in K_0$, apply the same cutoff construction to $H_+,H_-$
to obtain $p,h\in\Her(H_+)$ and $q,k\in\Her(H_-)$ for
Lemma~\ref{lem:finite-coupling}.  The contractions $p,q$ are
orthogonal.  The resulting
$U'=U'^*\in\mathcal D$ has a selected nonzero corner at every
point of $K_0$.  Put
\[
 x=H_+U'H_-\in\mathcal D.
\]
If $x_P=0$, then $(pU'q)_P=0$ for every
$p\in\Her(H_+)$ and $q\in\Her(H_-)$: approximate them in
norm by elements $H_+aH_+$ and $H_-bH_-$, and use
\[
 (H_+aH_+)U'(H_-bH_-)=H_+a\,x\,bH_-.
\]
This contradicts the selected nonzero corner at $P$.
Thus $x_P\ne0$ throughout $K_0$, while
$x^2=H_+U'(H_-H_+)U'H_-=0$.
All constructed elements lie in $\mathcal D$; spectral
projections have only been used to verify their supports.
\end{proof}

\subsection{Removing the amplification}

We have constructed a square-zero element on every prescribed
quasi-compact set where a hereditary subalgebra is nonzero.  We now
apply this result repeatedly on the same primitive norm slice.  The
resulting number of orthogonal equivalent copies can be made at least
the single amplification integer in \propF{}.

\begin{proof}[Proof of Theorem~\ref{thm:exact-infiniteness}]
Fix $a\in C_+$ and $r>0$.  The set $K_0=K(a,r)$ is quasi-compact by
Lemma~\ref{lem:primitive-compactness}.  If it is empty, every primitive
image of $(a-r)_+$ is zero, so $(a-r)_+=0$, and
$((a-r)_+)^{\oplus2}\precsim a$ is immediate.
Suppose that $K_0$ is nonempty.  Proposition~\ref{prop:square-zero}
applied to $\Her(a)$ gives $x\in\Her(a)$ with $x^2=0$ and
$x_P\ne0$ throughout $K_0$.  Set $z_1=x^*x$.  Its Cuntz-equivalent
positive $xx^*$ is orthogonal to it, since
$(x^*x)(xx^*)=x^*x^2x^*=0$.  Therefore
\[
 z_1^{\oplus2}\CuEq z_1\oplus xx^*
 \CuEq z_1+xx^*\precsim a,
 \qquad (z_1)_P\ne0\quad(P\in K_0).
\]
Repeat the construction inside $\Her(z_1)$, and then inside successive
hereditary subalgebras, always over the same $K_0$.  It yields
$z_j\in C_+$ nonzero throughout $K_0$ and
$z_{j+1}^{\oplus2}\precsim z_j$.  Finite direct sums and transitivity
give
\[
 z_j^{\oplus 2^j}\precsim a.
\]
Choose a finite $j\ge1$ with $2^j\ge n$.  Dropping extra summands gives
$z_j^{\oplus n}\precsim a$.  If $(z_j)_P=0$, then $P\notin K_0$, so
$\|a_P\|<r$ and $((a-r)_+)_P=0$.  The primitive ideal test in
Lemma~\ref{lem:ideals-cutoffs} therefore gives
$(a-r)_+\in\Id(z_j)$.  Condition \propF{} makes $z_j^{\oplus n}$ properly infinite because
$z_j$ is a scalar positive element.
Lemma~\ref{lem:ideal-comparison}, used with two copies on the left,
then implies
\[
 ((a-r)_+)^{\oplus2}\precsim z_j^{\oplus n}\precsim a.
\]
Only scalar positives $z_j\in C$ have been tested against \propF{}.
Finally, $\|a-(a-r)_+\|\le r$.  As $r\downarrow0$, norm-limit closure
of comparison with the fixed target $a$ gives
\begin{equation}\label{eq:individual-infiniteness}
 a\oplus a\precsim a\qquad(a\in C_+).
\end{equation}
This includes $a=0$; when $n=1$, it also follows directly from the
hypothesis.  Therefore $C$ satisfies \propP{}.
Theorem~\ref{thm:individual-to-strong} now gives strong pure
infiniteness, completing the proof.
\end{proof}

The removal of amplification also gives the comparison form needed
for the direct quantitative argument in Section~\ref{sec:quantitative}.
The following conclusion concerns the entire closed ideal, rather
than only one cutoff or primitive norm slice.

\begin{corollary}[Whole-ideal comparison]\label{cor:ideal-comparison}
Let $C$ be a $C^*$-algebra satisfying \propP{}.  For every $h\in C_+$,
$c\in\Id(h)_+$, and integer $k\ge1$,
\[
 c^{\oplus k}\precsim h.
\]
If $C$ is exact, Proposition~\ref{prop:perturbation} may be applied
with multiplicity $n=1$, with the same constant $\sigma$ and all its
stated options.
\end{corollary}

\begin{proof}
Property \propP{} makes $h$ properly infinite, so
Lemma~\ref{lem:ideal-comparison} applies with multiplicity one.
If $h=0$, the ideal is zero and the assertion is immediate.
For exact $C$, property \propP{} is precisely the amplification
hypothesis of Proposition~\ref{prop:perturbation} with $n=1$ for every
scalar positive element.  Its constant was fixed independently of $n$,
so neither the constant nor the hereditary and selfadjoint options change.
\end{proof}
